\textbf{Correspondence status: not reviewed. The shared module and section companion compile with the pinned Lean and Mathlib versions. The axiom report for all section and shared-module declarations contains only the standard foundational axioms \texttt{propext}, \texttt{Classical.choice}, and \texttt{Quot.sound}; these machine checks do not replace Steven's review.}

The formalization uses Mathlib's actual matrix multiplication, trace, characteristic polynomial, polynomial roots, bases, and endomorphism trace. The shared results are in \texttt{TraceFormulas.LinearAlgebra.MatrixTrace}; this section specializes them to complex matrices indexed by $\operatorname{Fin}(n)$. Lean's indices $0,\ldots,n-1$ correspond to the book's indices $1,\ldots,n$. The statements allow $n=0$, and rectangular cyclicity allows either dimension to be zero.

For the \hyperref[linear-algebra.trace-eigenvalue-sum]{trace--eigenvalue theorem}, the right side is the sum of the characteristic polynomial's root multiset. A multiset retains repeated roots, so this represents algebraic multiplicity rather than distinct eigenvalues or eigenspace dimensions. The companion also proves that this multiset has cardinality $n$. The statement has no diagonalizability assumption. Mathlib's proof uses the same coefficient comparison as the prose, with a separate empty-index case. The complex algebraic-closure instance is Mathlib's proved fundamental theorem of algebra, not a project axiom.

The \hyperref[linear-algebra.trace-cyclic]{cyclicity theorem} uses two possibly different matrix sizes. The \hyperref[linear-algebra.trace-similarity]{similarity theorem} assumes \texttt{IsUnit S}, the precise assertion that the matrix $S$ is invertible, and uses Mathlib's matrix inverse in $S^{-1}AS$. The trace definition, additivity, scalar multiplication, identity-matrix trace, and empty trace are also exposed in the companion. Characteristic roots are related to nonzero eigenvectors through the determinant of $zI-A$ and Mathlib's singular-matrix kernel theorem.

For the \hyperref[linear-algebra.endomorphism-trace]{endomorphism trace}, a finite basis is an explicit input. The shared module proves that Mathlib's \texttt{LinearMap.trace} agrees with the matrix trace in that basis, and that the matrices in any two finite bases have equal traces. Although the library defines an endomorphism trace more generally, these statements require the finite basis used in the manuscript.

The opening example checks the trace, both displayed eigenvector equations, and that the vectors are nonzero. The rectangular example checks both products and both traces. The Jordan example checks its characteristic polynomial, the root multiset $\{2,2\}$, trace $4$, its eigenspace as the span of $(1,0)^{\mathsf{T}}$, and the eigenspace dimension $1$. A further theorem rules out every basis of eigenvectors, so the claim of non-diagonalizability is proved explicitly.

The checked library sources include \texttt{Matrix.Charpoly.Eigs}, \texttt{Matrix.Trace}, \texttt{Matrix.NonsingularInverse}, \texttt{Matrix.ToLinearEquiv}, and \texttt{LinearAlgebra.Trace}. They are used at the project's exact dependency revision. Source hashes, build commands, and axiom evidence belong to the formalization audit; the source links and website build do not themselves establish these facts. Steven's review of the meanings, conventions, and correspondence is still pending, and no video approval follows from these machine checks.
